\documentclass{article}

\usepackage{fancyhdr}
\usepackage{extramarks}
\usepackage{amsmath}
\usepackage{amsthm}
\usepackage{amsfonts}
\usepackage{tikz}
\usepackage[plain]{algorithm}
\usepackage{algpseudocode}
\usepackage{fullpage, url, hyperref}


\theoremstyle{plain}
\newtheorem{theorem}{Theorem}
\newtheorem{lemma}{Lemma}
\newtheorem{definition}{Definition}

\newcommand{\I}{\mathcal{I}}
\newcommand{\F}{\mathbb{F}}
\newcommand{\E}{\mathbb{E}}
\newcommand{\R}{\mathbb{R}}
\renewcommand{\P}{\mathbb{P}}
\usetikzlibrary{automata,positioning}

%
% Basic Document Settings
%

\topmargin=-0.45in
\evensidemargin=-.50in
\oddsidemargin=-.5in
\textwidth=7in
\textheight=9.5in
\headsep=0.25in

\linespread{1.1}

\pagestyle{fancy}
\lhead{\hmwkAuthorName}
\chead{\hmwkClass\ (\hmwkClassInstructor\ \hmwkClassTime): \hmwkTitle}
\rhead{\firstxmark}
\lfoot{\lastxmark}
\cfoot{\thepage}

\newcommand{\w}[2]{w_{#1}^{(#2)}}
\newcommand{\loss}[2]{\vec{\ell}_{#2}^{(#1)}}
\newcommand{\losst}[1]{\vec{\ell}^{(#1)}}
  \newcommand{\pt}[1]{\vec{p}^{(#1)}}
\renewcommand\headrulewidth{0.4pt}
\renewcommand\footrulewidth{0.4pt}

\setlength\parindent{0pt}


% 
% Create Problem Sections
%


\newcommand{\enterProblemHeader}[1]{
    \nobreak\extramarks{}{}\nobreak{}
    \nobreak\extramarks{}\nobreak{}
}

\newcommand{\exitProblemHeader}[1]{
    \nobreak\extramarks{}\nobreak{}
    \stepcounter{#1}
    \nobreak\extramarks{}\nobreak{}
  }

  
\newcommand{\enterExerciseHeader}[1]{
    \nobreak\extramarks{}{}\nobreak{}
    \nobreak\extramarks{}\nobreak{}
}

\newcommand{\exitExerciseHeader}[1]{
    \nobreak\extramarks{}\nobreak{}
    \stepcounter{#1}
    \nobreak\extramarks{}\nobreak{}
}

\setcounter{secnumdepth}{0}
\newcounter{partCounter}
\newcounter{homeworkProblemCounter}
\setcounter{homeworkProblemCounter}{1}

\newcounter{homeworkExerciseCounter}
\setcounter{homeworkExerciseCounter}{1}
\nobreak\extramarks{}\nobreak{}
\nobreak\extramarks{}\nobreak{}

%
% Homework Problem Environment
%
% This environment takes an optional argument. When given, it will adjust the
% problem counter. This is useful for when the problems given for your
% assignment aren't sequential. See the last 3 problems of this template for an
% example.
%
\newenvironment{homeworkProblem}[1][-1]{
%    \ifnum#1>0
   %     \setcounter{homeworkProblemCounter}{#1}
   % \fi
    \section{Problem \arabic{homeworkProblemCounter}}
    \setcounter{partCounter}{1}
    \enterProblemHeader{homeworkProblemCounter}
}{
    \exitProblemHeader{homeworkProblemCounter}
}

  \newenvironment{homeworkExercise}[1][-1]{
    \ifnum#1>0
        \setcounter{homeworkExerciseCounter}{#1}
    \fi
    \section{Exercise \arabic{homeworkExerciseCounter}}
    \setcounter{partCounter}{1}
    \enterExerciseHeader{homeworkExerciseCounter}
}{
    \exitExerciseHeader{homeworkExerciseCounter}
}

%
% Homework Details
%   - Title
%   - Due date
%   - Class
%   - Section/Time
%   - Instructor
%   - Author
%

\newcommand{\hmwkTitle}{Homework\ \#4}
\newcommand{\hmwkAssignedDate}{March 26}
\newcommand{\hmwkDueDate}{April 9, midnight}
\newcommand{\hmwkClass}{Design and Analysis of Algorithms, CS 6550}
\newcommand{\hmwkClassTime}{}%Section A}
\newcommand{\hmwkClassInstructor}{Professor Jamie Morgenstern}
\newcommand{\hmwkAuthorName}{}%\textbf{Josh Davis} \and \textbf{Davis Josh}}

%
% Title Page
%

\title{    \vspace{2in}
    \textmd{\textbf{\hmwkClass:\ \hmwkTitle}}\\
    \normalsize\vspace{0.1in}\small{Due\ on\ \hmwkDueDate\ at 3:10pm}\\
    \vspace{0.1in}\large{\textit{\hmwkClassInstructor\ \hmwkClassTime}}
    \vspace{3in}
}

\author{\hmwkAuthorName}
\date{}

\renewcommand{\part}[1]{\textbf{\large Part \Alph{partCounter}}\stepcounter{partCounter}\\}

%
% Various Helper Commands
%

% Useful for algorithms
\newcommand{\alg}[1]{\textsc{\bfseries \footnotesize #1}}

% For derivatives
\newcommand{\deriv}[1]{\frac{\mathrm{d}}{\mathrm{d}x} (#1)}

% For partial derivatives
\newcommand{\pderiv}[2]{\frac{\partial}{\partial #1} (#2)}

% Integral dx
\newcommand{\dx}{\mathrm{d}x}

% Alias for the Solution section header
\newcommand{\solution}{\textbf{\large Solution}}

% Probability commands: Expectation, Variance, Covariance, Bias
\newcommand{\Var}{\mathrm{Var}}
\newcommand{\Cov}{\mathrm{Cov}}
\newcommand{\Bias}{\mathrm{Bias}}
\usepackage{enumitem}
\newcommand\litem[1]{\item{\bfseries#1.\space}}


\begin{document}

% \maketitle

\pagebreak

{\bf Homework Out: \hmwkAssignedDate}\\
{\bf Due Date:  \hmwkDueDate}\\

The HW contains some exercises (fairly simple problems to check you
are on board with the concepts; don’t submit your solutions), and
problems (for which you should submit your solutions, and which will
be graded). Some problems have sub-parts that are exercises.  For this
problem set, it’s OK to work with others. (Groups of 2, maybe 3 max.)
That being said, please think about the problems yourself before
talking to others. Please cite all sources you use, and people you
work with. The expectation is that you try and solve these problems
yourself, rather than looking online explicitly for
answers. Submissions due at beginning of class on the due date.
Please check the Piazza for details on submitting your \emph{LaTeXed}
solutions.

\newcommand{\poly}{\textrm{poly}}
\paragraph{Exercises}
\begin{enumerate}

  \litem{Lots of Flows} Suppose you wanted to find an approximate
  solution to the following “multicommodity” flow problem: given a
  digraph $G = (V, E)$ with unit arc capacities, send $F_i$ flow from
  node $s_i$ to node $t_i$ in the graph, for all $i \in [k]$. You
  should imagine that the flow from $s_i$ to $t_i$ is of commodity $i$
  (e.g., oil, water, sand...) which are all distinct.
  \begin{enumerate}
\item     Suppose $\P_i$
is the set of all paths from $s_i$ to $t_i$: show that the following LP captures the
problem we are trying to solve. The variables are $f_P$ , one for each path in $\cup_i\P_i$.

\begin{align}
  \left.\begin{array} { ll }
          { \sum _ { P \in P _ { i } } f _ { P } = F _ { i } } & { \forall i \in [ k ] } \\
          { \Sigma _ { i } \sum _ { P \in P _ { i } e \in P } f _ { P } \leq 1} & { \forall e \in E }
        \end{array} \right.
\end{align}

\item Define an appropriate “easy” polytope $K$ for this problem.
  \item  Given weights $q \in \Delta^m$, how would you solve the oracle for this problem? Show you can
find a flow that satisfies the demands, but uses at most $(1 + \epsilon)$ capacity on each edge in time $\frac{O(k(m+n \ln n ))}{\poly(\epsilon)} \cdot \left(\sum_i F_i\right) = \frac{O(km(m+n\ln n))}{\poly(\epsilon)}$. 
\end{enumerate}
\litem{Strength in Convexity} A function $f : \R^n \to \R$ is called $\ell$-strongly convex if
\begin{align*}
f ( y ) \geq f ( x ) + \langle \nabla f ( x ) ,y - x ) + \frac {ell } { 2} | | y - x | | ^ { 2}
\end{align*}
%
I.e., if the function is not just convex, but “locally it grows at least as fast as a quadratic”.
Modify the basic gradient descent analysis to show that using the same update rule $x_{t+1}\gets
x_t − \eta_t \nabla f(x_t) $with suitably chosen $\eta_t$, then we can find $\hat{x}\in \R^n$
such that
\begin{align*}
f ( \hat { x } ) - f ( x ^ { * } ) \leq O ( \frac { G ^ { 2} \log T } { \ell \cdot T } ).
\end{align*}
Again, assume that $||\nabla f(x)|| \leq G$. Note due to the
assumption of strong convexity, we got better convergence: the
dpendence on $T$ is better, there's no dependence on
$D = ||x_0 - x^*||$. Show that this analysis also works in the online
case if each function is strongly convex. Bonus: try and remove the
$\log T$ term in the offline case. Why does this not work in the
online case?
\litem{Divergent Views.} Given two discrete probability distributions $p, q$ defined over a universe of $N$ elements, the Kullback-Liebler divergence between the two is defined as
\[KL(p || q) = \sum_{i = 1}^N p_i \log_2 \frac{p_q}{q_i}.\]
%
Show the following:
\begin{enumerate}
\item Give examples where $KL(p || q)\neq KL(q || p)$ where both divergences are finite.
\item Show $KL(p || q) \geq 0$ and $KL(p || q) = 0$ only if $p =q $.
\item If $U_N$ is the uniform distribution over $N$ elements, then $KL(p || U_n) \leq \log_2 N - H(p)$ where $H(p)$ is the Shannon entropy of $p$.
\item Show that the Bregman divergence $D_h(p || q) = KL(p||q)$ when $h(x)= \sum_i x_i \log_2 x_i$ is the negative entropy function.
\end{enumerate}
And here are some useful facts about Bregman divergences. Recall that $D_h(x||y) = h(x) - h(y) - \langle \nabla h(y), x-y\rangle$ for a strinctly convex $h$.
\begin{enumerate}
  \item[(e)] Given points $x_1, _2 \ldots, x_n$, show that the unique point $c$ that minimizes the average distance $\frac{1}{n} \sum_i D_h(x_i || c)$ is the center of mass $c = \frac{1}{n} \sum_i x_i$.
\end{enumerate}
\litem{That's the Norm.} In lecture, we defined a differentiable convex function $f : K \to \R$ to be $L$ Lipschitz with respect to $|| \cdot ||$ if
\[\frac{| f(x) - f(y)|}{||x - y||} \leq L.\]
Show that this is equivalent to $|| \nabla f(x)||_* \leq L$ for all $x\in K$.

Similarly, show that $f$ being $\alpha$-strongly-convex with respect
to $|| \cdot ||$ is equivalent to
$|| \nabla f(x) - \nabla f(y)||_* leq \beta || x - y||$.
\end{enumerate}
\paragraph{Problems}


\begin{enumerate}
  \litem{Hmm, that's odd...} To solve the max-weight perfect matching
  problem, we need to optimize over the perfect matching polytope. In
  turn, this requires that we find a separation oracle for the odd cut
  constraints. I.e., given $x\in \R^{|E|} $ we need to find
  $S \subseteq V$ such that $|S|$ is odd, and $x (\partial S)$ is
  minimized, where
  $x ( \delta(S)) = \sum_{i \in S, j \not \in S} x_{ij}$. Then,
  comparing this min-odd-cut value to $1$, we can find a violated
  constraint if one exists. Assume $|V| $ is even, otherwise the LP is
  infeasible.
  \begin{enumerate}
  \item A function $f : 2^V \to \R$ is \emph{submodular} if for all $A, B \subseteq V$, we have
    \[f(A) + f(B) \geq f(A \cup B) + f (A \cap B).\]
    %
    Show that $f(A) = x (\partial(A))$ is submodular. Observe $f$ is symmetric: $f(A) = f(V \setminus A)$.
  \item If $(C, \bar{C})$ is the min cut, i.e. if $x (\partial(C))$ is
    the least among all non-empty cuts, and $|C|$ is even, then show
    that there exists a min-odd cut either contained within $C$ or
    $\bar{C}$.
    \item Give an algorithm to find a min-odd-cut in polynomial time. 
    \end{enumerate}

    \litem{Zero-Sum Games using LP Duality.} Recall the zero-sum game setup: we have some matrix $M \in \R^{m \times n|}$; if the row player plays $x \in \Delta^m$ and the column player plays $y \in \Delta^n$, the payoff for the row player is $x^T M y$.

    If we define $C(x) = \min_{y \in \Delta^n} x^T M y$ and $R(y) = \max_{x \in \Delta^m} x^T M y$, the minimax theorem proves that (a) for all $x,y, C(x) \leq R(y)$, and (b) there exist $x^*, y^*$ such that $C(x) = R(y)$.
    \begin{enumerate}
    \item Show an LP to compute $\max_x C(x)$, the optimal strategy for the row player. Hint: be careful, the definition of $C$ has a  minimum in its definition, looking to find $\max_x \min_y x^T M y$, which in this form is not a linear program.
    \item Show an LP to compute $\min_y R(y)$, the optimal strategy for the column player.

    \item Show that you can in fact find LPs for the above such that the optimal solution for the dual for the first LP is a solution to the second LP.
    \item (Do not submit.) Use weak duality to infer the first part of
      the minimax theorem, and strong duality to infer the second
      part.
      \end{enumerate}
      \litem{Capacitated Max-Flow and Width Reduction.} Consider the directed $s-t$ max flow problem: $K =  \{f | f_P \geq 0, \sum_{P \in \mathcal{P}} f_P = F\}$, with constraints
      \[f_e / c_e \leq 1 \hspace{.5in} \forall e\in E,\]
      %
      where $f_e = \sum_{ P : e \in P} f_P$. In lecture we considered $c_e = 1$, now we consider the general case. Given the weights $p \in \Delta^m$ given by Hedge, the ``average'' constraint looks like
      \[\sum_e p_e (f_e/c_e) = \sum_e f_e p_e/c_e \leq 1.\]
      \begin{enumerate}
      \item {Do not submit} Suppose the oracle sends the entire $F$ units of flow along a shortest path w.r.t. $p_e / c_e$. show that there are capacitated networks, and possible $p\in \Delta^m$, where all $F$ flow is routed along a path using the least-capacity edge. Hence the width of this oracle is at least $(F/c_{\min})$.

        Note that $F$ may be as large as $\Omega(m c_{\max})$ so with
        general capacities this ratio could be much larger than
        $m$. Now, we'll try the idea of bumping all of the edge
        weights up a bit to reduce the width of the problem at the
        expense of some approximation in our solution.  For the
        remainder of the problem, assume $\epsilon \leq
        \frac{1}{10}$. You may assume the instance is feasible (one
        can send a flow of value $F$ while satisfying all capacity
        constraints).

      \item Set the weights to be $w_e = p_e + \epsilon/m$. Compute the shortest path $P^*$ w.r.t. edge lengths $\frac{w_e}{c_e}$, and send all $F$ flow along it. Show
        \[F \cdot \sum_{e \in P^*} \frac{p_e}{c_e} \leq \min_{f\in K} \sum_{e\in E} \frac{f_e}{c_e} w_e \leq 1 + \epsilon.\]
      \item Show that
        \[ \max_{e\in P^*} F/c_e \leq O(m/\epsilon)\]
        and so the width of the oracle is $O(m/\epsilon)$.
      \item (Do not submit.) Using this oracle and the MW algorithm,
        give an $\tilde{O}(m^2/\epsilon^3)$-time
        $(1+\epsilon)$-approximate max-flow algorithm for capacitated
        graphs.
      \item Bonus: use these ideas to get an algorithm for the
        multi-commodity case from exercise \# 1 that works for
        capacitated graphs, but whose runtime does not depend on the
        magnitude of those capacities.

        Note: This idea was used in the Christiano et al. paper,
        combined with electrical flows, to bring the width down to
        $O(\sqrt{m/\epsilon})$.
      \end{enumerate}
      \litem{Solving a Linear System.} Given a positive-definite $A \in \R^{n \times n}$ with eigenvalues $0 < \lambda_1 \leq \lambda_2 \leq \ldots \leq \lambda_n$, the condition number of $A$ is $\kappa = \frac{\lambda_n}{\lambda_1}$. Given a vector $b\in \R^n$, the goal is to find a near-solution to the linear system $Ax = b$. Consider the function $f(x) = \frac{1}{2} x ^T A X - bx$.
      \begin{enumerate}
      \item (Do not submit) Show that $f$ is convex, and $\nabla f(x) = Ax - b$. Hence infer that the minimizer $x^*$ of $f$ satisfies $A x^* = b$.  Also show that $f$ is $\lambda_1$-strongly convex and $\lambda_n$-smooth.
      \item Show that gradient descent on $f$ starting at some $x_0 \in \R^n$ guarantees
        \[ || x_t - x^*||_2 \leq \max\{|\mu_1|, |\mu_n |\}^t \cdot || x_0 - x^*||_2\]
        where $\mu_1 \leq \ldots \leq \mu_n$ are the eigenvalues of $(I - \eta A)$.
      \item Show that $||x_t - x^*||_2 \leq \epsilon \cdot || x_0 - x^*||_2 $ after $O(\kappa \log \frac{1}{\epsilon})$ iterations.
      \item (Do not submit.) Define $||x||_A = \sqrt{x^T A x}$. Show
        that $|| x_t - x^*||_A \leq \epsilon || x_0- x^*||_A$ after
        $O(\kappa \log\frac{\kappa}{\epsilon})$ iterations.
      \end{enumerate}
      \litem{Gradient Descent, meet Linear Optimization.} When we did
      constrained GD over a convex $K$, we took a step along the
      negative gradient and projected back to $K$. Here is a different
      approach using LPs. Give $x_t \in K$ the next iterate is
      \begin{align*}
        y_t & \gets \arg\min_{y \in K}\{(\nabla f  (x_t))^T y\} \\
        x_t & \gets (1-\eta_t) x_t + \eta_t y_t.              
      \end{align*}
      The minimizer in the first step can be found using linear optimization over $k$, which may be much simpler than general convex optimization.
      \begin{enumerate}
      \item Assume that the function $f$ is $\beta$-smooth w.r.t. some norm $|| \cot ||$, and $R = \max_{a,b \in K} || a - b||$ in the same norm, and $\eta_t = \frac{2}{t+1}$. Show that for all $t\geq 0$:
        \[f(x_{t+1}) - f(x_t) \leq \eta_t \left(f(x^*) - f(x_t)\right) + \frac{\beta}{2} \eta^2_t R^2\]
        \item Use induction to show that $f(x_t) - f(x^*) \leq \frac{2\beta R^2}{t+1}$.
        \item (Do not submit.) Suppose $K$ is a polytope with $n$
          non-negativity and $m$ other constraints. Show that we can
          choose $y_t$ to be a vertex of $K$ with at most $m$ non-zero
          coordinates. So, if we start with $x_0 - \bar{o}$, $x_t $
          will have support over at most $mt$ coordinates.  Show the
          previous problems objective guarantee for this $x_t$.
        \end{enumerate}

          Now we use this algorithmic setup to solve the ``topic
          model'' problem. Consider a topic matrix
          $A\in \R^{m \times n}$, with columns being associated with
          topics and rows being associated with words. Let $y\in \R^m$
          be a document, and we'd like to find a set of topics $x$
          such that $y$ was likely generated by those topics:          \[ y \textrm { is close to } Ax \textrm{    and   } x \textrm{ is sparse.}\]
          %
          Since $||x||_0$  minimization is hard, we often use $||x ||_1$ minimization instead:
          \[\min_{x\in \R^n} || y - A_x||_2^2 \textrm{ s.t. } ||x||_1 = 1.\]
          In this setting $n \gg m$, many many many topics but many fewer words, and we'd ideally like to avoid running in $\poly(n)$ time. Suppose we have an oracle that given $w\in \R^m$ outputs $\arg\min_{j\in[n]} \langle w, A_j\rangle$ in $Z \geq m$ time.
            \begin{enumerate}
          \item[(d)] Show that each iterate $x_t$ can be compute in $O(Zt + t^2)$ time, by keeping track of the nonzero entries in $x_t$.
          \item[ (e)] Show that if each column length $||A_j||_2 \leq L$ then $f(x) = || y - Ax||^2_2$ is $2L^2$-smooth w.r.t. the $\ell_1$ norm. You may use exercise 4 without proof. Observe that the $\ell_1$-diameter $R$ of the polytope $K = \{x | ||x||_1 = 1\}$ is $1$.
            
            So, we get an $\epsilon$-approximate solution in time $O(ZL^2/\epsilon + L^4/\epsilon^2)$.
          \end{enumerate}
\end{enumerate}



\end{document}
